\documentclass[dvipdfmx]{jlreq}

\title{マイポリシー}
\author{船原滉樹}
\date{\today}

\begin{document}
\maketitle

\section{はじめに}
    私は、世界の本質を探究し続ける。
    常に広く、深く探求する。
    将来的にはAIの基盤技術を開発・運用することを目指す。
    さらに宇宙分野への挑戦も視野に入れる。
    長期的には、社会や世界に強い影響力を持つ存在になることを目標とする。
    本書は、そのための思想と行動指針をまとめたものである。

\section{チャレンジ}
    失敗を恐れず、常に新しいことに挑戦する。

\section{振り返り}
    無理のない範囲で、継続的に振り返りを行う。
    毎月上旬に、先月の振り返りを行う。方法は、先月のSNSの投稿などを見返す。
    毎月下旬に、今月の定量評価を行う。方法は、Awarefyのストレスチェックで過去一カ月間のストレスレベルを測る。

\section{生命維持}
    \subsection{生理的欲求}
        生命維持のために、生理的欲求を満たすことを優先する。特に、睡眠時間を確保することを優先する。

    \subsection{自死対策}
        どうしても死にたいと思わされてしまうときもあるかもしれない。
        そういう状況下でも命を繋いでいくためには、以下の方法を試す。
        \begin{itemize}
            \item 今は死にたいと思っていることを認める。
            \item 信頼できる人、またはAIに相談する。
            \item 精神科で薬物療法も検討する。
            \item マインドフルネスを実践する。
            \item 物理的な刺激を与える。

        \end{itemize}
    \subsection{マイルールを破る}
        生命維持のためにマイルールを破ることは立派な選択である。


\section{目的意識}
    できる限り、目的を持って行動する。
    ただし、目的は必ずしも一つとは限らない。

\section{健康管理}
    \subsection{睡眠}
    \subsubsection{睡眠時間}
        毎日7時間以上の睡眠を確保する。8時間以上が望ましい。
    \subsection{メンタルヘルス}
        私は、メンタルの安定を保つ。
        無理をしない。

\section{情報セキュリティ}
    \subsection{スマホのパスワード}
        スマホのパスワードは、8桁以上の英数字を含むものとする。12桁以上が望ましい。
    \subsection{パソコンのアプリのインストール}
        自分のパソコンにアプリをインストールする際には、以下の点に注意する。その際、AIを活用して、アプリの安全性を確認してもいい。
    \begin{itemize}
        \item アプリの提供元が信頼できるかどうかを確認する。
        \item アプリの権限を確認する。
        \item 偽のリンクに注意する。
    \end{itemize}

\section{思考整理}
    \subsection{マインドマップ}
        \subsubsection{紙に書き出す}
\section{学習方法}
    \subsection{アプトプット}
        \subsubsection{目的}
        \begin{itemize}
            \item 学習内容の定着のため。
        \end{itemize}



\section{優先順位}
    迷ったときは、以下のリストを参考にする。
    \begin{enumerate}
        \item 10年後の自分にとって重要なこと
        \item 1年後の自分にとって重要なこと
        \item 1ヶ月後の自分にとって重要なこと
        \item 1週間後の自分にとって重要なこと
        \item 今日の自分にとって重要なこと
        \item 安全欲求を大事にする
    \end{enumerate}

\section{人間関係}
    \subsection{コミュニケーションスタイル}
        \subsubsection{傾聴ルール}
        \begin{itemize}
            \item 相手の感情を理解しようとする。
        \end{itemize}

\section{PDCAサイクル}
    一年間を通してや、一カ月を通して、一つの活動を通して等、PDCAサイクルを回す。
    \subsection{根本原因を特定方法}
        5回「なぜ」を繰り返す。
        マインドマップを作製する。

\section{生成AI、AIエージェント}
    効率化を図るために、AIエージェントを活用する。
    \subsection{注意点}
    AIエージェントを活用する際には、以下の点に注意する。
    \begin{itemize}
        \item 開発において、システムの仕組みを理解する。
    \end{itemize}



\section{行動の評価方法}
    \subsection{加点方式}
        行動活性化のために、加点方式を採用する。また、メンタル的にも辛くないため、加点方式を採用する。
    \subsection{減点方式}
        精神的・身体的にエネルギーが充分に溜まっているのに行動が停滞しているときに、刺激を与える目的で、減点方式を採用する。スピード感がモチベーションに繋がる場合もあるため。
    \subsection{定量的な評価方法}
        物事を評価する際には、定量的な評価方法を採用する。

\section{進捗度算出}
    \subsection{50\%-50\%法}
        次の3段階で進捗度を算出する。
        \begin{table}[ht]
            \centering
            \begin{tabular}{|c|c|}
                \hline
                状態 & 進捗度 \\
                \hline
                未着手 & 0\% \\
                \hline
                着手済み & 50\% \\
                \hline
                完了 & 100\% \\
                \hline
            \end{tabular}
            \caption{50\%-50\%法}
        \end{table}
        プロジェクト全体の進捗度は、各タスクの進捗度の平均値で算出する。

        進捗管理は「ゆるたす」や「Google Calendar」等を使う

\section{自分の考えや価値観に矛盾が生じたとき}
    状況によって使い分けをすることが大事だと思う。使い分ける方法については、これから考える。

\section{マイルールの破り方}
    生命維持や成長の目的で、マイルールを破ることは立派な選択である。
    ただし、ルールを破る際には、以下の点に注意する。
    \begin{itemize}
        \item ルールを破る理由を明確にする。
        \item ルールを破ることで得られるメリットとデメリットを評価する。
        \item ルールを破ることで、他人に迷惑をかけないようにする。
        \item ルールを破ることで、ルールの効力が下がってしまう場合があることを理解する。
    \end{itemize}
    

\section{マイルールの変更}
    マイルールの変更はいつでも可能である。ルールを破った時や、ルールに対して不満を感じたときが変更のタイミングだと思う。ただし、変更する際には、以下の手順で行う。
    \begin{enumerate}
        \item ルールに対する不満や改善点を明確にする。
        \item 変更案を考える。
        \item 変更案を実行する前に、メリットとデメリットを評価する。
        \item 変更案を実行する。
    \end{enumerate}


\end{document}
